% Nouveau EGA V (cote 25): archive cote 25, site page 176 (PDF page 177).
% typescript. Draft compared: none.
% AI-transcribed by Claude Opus 5.5 (claude-opus-5-5, adaptive thinking, effort high), Message Batches API, batch ega25b (msgbatch_01A7n5YuXubUPNK1zsepi8RG), 2026-11-13; not yet checked by eye.
\scanpage{25/176}
\mspage{\hand{8}}

coéfficients entiers en les $a_i,b_{ij}$ . Alors le ``contenu'' de ce polynôme
est égal à 2 (i.e. le pgcd de ses coéfficients est 2).
\inserted[typed]{multiple de 2.}
\xout{Cela} Prouvons d'abord que ce contenu est \xout{de la forme 2\deleted{$^r$}\illegible{1}}
En effet, cela signifie aussi que \deleted{\xout{\illegible{1}}} lorsqu'on spécialise ce polynôme
en \xout{la}un corps $Z/\over{p}{2}Z$\hand{)} \deleted{avec $p\neq 2$}, il \xout{n'}est pas identiquement nul, ce qui \xout{xxx}
résulte du fait qu'une forme quadratique de degré imp\over{q}{a}ir sur un corps de
car 2 est tjrs dégénérée, i.e. à un discriminant nul. Pour prouver
que le contenu \inserted{c} est exactement 2, il suffit de fai\unsure{r}e le calcul du discri-%
minant \inserted{de la forme} standart de degré n, d'une part elle doit être multiple de c,
d'autre part le calcul donne 2, (car c'est 1 pour le degré pair, et 2 pour
le degré 1 ...).

\xout{\unsure{On en conclut immédiatement}} Il y a lieu alors d'introduire le polynôme
\[ \overset{\hand{\sim}}{d}\xout{*}(Q) = \tfrac{1}{2}d(Q) \]
en les coéfficients de Q, qui est encore à coéfficients entiers, donc
prend une valeur bien déterminée quand on spécialise les $a_i,b_{ij}$ dans un
anneau A quelconque, i.e. quand on prend une forme quadratique \deleted{\xout{Q}}\inserted{q}de degr\unsure{é}
n à coéfficients dans A. On appela $\overset{\hand{\sim}}{d}\xout{*}(Q)$ le``\uline{polynôme} \unsure{\hand{``}}\uline{discrimihant}
\uline{corrigé de la forme quadratique indétérminée Q à n variables}'' , et sa
valeur relativement aux coéfficients d\xout{e}'une forme quadratique q à coeff
dans un anneau A quelconque sera appelé le ``\uline{discriminant corrigé de la}
\uline{forme quadratique} q''. Plus généralement, on déduit de façon essentielle
ment triviale de 16.5. l'énoncé suivant:

\uline{Proposition} 16.6. On peut, \inserted{et} d'une seule manière, associéer à tout\xout{e}
préschéma S, muni d'un Module loc libre \xout{d}e rang fini \uline{E} et d'une forme \inserted{qu.}
$Q \in \unsure{\Gamma}\,\underline{\mathrm{Sym}}^2(E)$, associer une section $\overset{\hand{\sim}}{d}(Q)$ de dét$(E)^{\otimes 2}$ , de façon à satis-%
faire aux conditions suivantes:

(\xout{1}a) Compatibilité avec changement de base\hand{, et fonctorialité pour isomorphismes en \uline{E}.}

(b) Si E est partout de rang pair, on a $\overset{\hand{\sim}}{d}(Q) =$ d(Q) (ou d(Q)
désigne le discriminant de Q).

\mspage{\hand{176}}
